% Appendix A. Mathematical background: introduction, index, Part A (categories) and Part B (monoidal categories,
% string diagrams and lenses). Parts C and D follow in sections/A2-background.tex.
% Source: site/sections/13-appendix.html (#app-how, #part-a, #part-b), converted in full; entry numbers A.1--B.17 kept.
% Statements about numbered results follow the public statements in theory/propositions.json.
\providecommand{\fbAusedin}[1]{\par\smallskip\noindent{\small\color{inkgray}{\sffamily\bfseries\color{tide}Used in.}\enspace #1\par}\medskip}

\section{Mathematical background}\label{app:background}

This appendix is for readers who know linear algebra, multivariable calculus, basic probability and the PEFT literature,
and who have not studied category theory, Lie groups or real algebraic geometry. It skips what ML readers already use,
such as the SVD, rank, orthogonal matrices, eigenvalues, Jacobians, gradients, Adam, LoRA and the transformer block. Each
of the 93 entries gives a short definition or statement with its hypotheses, a small example from fine-tuning, the places
in the paper that use it, and usually a textbook for further reading. Definitions carry an ochre rule and standard facts
quoted from the literature a teal one; the one remark, \bgref{A.20}, is set like a definition.

Part A covers categories, slices and base change, Part B monoidal categories, string diagrams and lenses, Part C
manifolds, cones and real algebraic geometry, and Part D groups, symmetry and conservation laws. The parts can be read in
any order, since each entry refers to the entries it depends on at their first use. The exposés cite an entry at the
first use of its term, as in \bgref{A.8}, and each entry ends with the results that rely on it.

A few symbols from the exposés recur in the examples. A method $M=(Q,q_0,\rho)$ is a smooth map $\rho:Q\to\Theta$ from
trainable parameters to weights with $\rho(q_0)=\theta_0$ (\bgref{A.8}), and $\lvert M\rvert=\dim Q$ is its parameter
budget. The expressive dimension $d(M)$, the first-order image $S_1$ and the first-order deficit are defined in
\bgref{C.3}; the base point is \emph{regular} for $M$ when the deficit is $0$ and an \emph{apex} when it is positive. LoRA
\citep{hu2021lora} is $\rho(B,A)=W_0+sBA$ with $B\in\R^{m\times r}$ and $A\in\R^{r\times n}$, and its full-rank
pointings fall into the strata $\mathrm T_A$, $\mathrm T_B$ and $\mathrm S$ of \bgref{C.19}. \emph{Block-scalar rates}
give each factor one learning rate, $\eta_B$ for $B$ and $\eta_A$ for $A$ (\bgref{D.23}).

Table~\ref{tab:bg-index} lists every entry with its area and the first places it is used.

{\footnotesize\setlength{\tabcolsep}{4pt}\renewcommand{\arraystretch}{1.08}
\begin{longtable}{@{}>{\raggedright\arraybackslash}p{0.8cm}>{\raggedright\arraybackslash}p{4.9cm}>{\raggedright\arraybackslash}p{2.35cm}>{\raggedright\arraybackslash}p{5.95cm}@{}}
\caption{Index of the background entries, with the area of each and the first places in the paper that use it.}
\label{tab:bg-index}\\
\toprule\headrow\textbf{No.} & \textbf{Entry} & \textbf{Area} & \textbf{First places used}\\\midrule
\endfirsthead
\toprule\headrow\textbf{No.} & \textbf{Entry} & \textbf{Area} & \textbf{First places used}\\\midrule
\endhead
\midrule\multicolumn{4}{r@{}}{\textit{continued on the next page}}\\
\endfoot
\bottomrule
\endlastfoot
\hyperref[bg:A.1]{A.1} & category & Categories & Table~\ref{tab:notation}, \thmref{Def.}{I.2}\\
\hyperref[bg:A.2]{A.2} & commutative diagram & Categories & Figure~\ref{fig:interval}, \thmref{Def.}{I.2}\\
\hyperref[bg:A.3]{A.3} & isomorphism & Categories & \thmref{Obs.}{I.6}, \thmref{Obs.}{I.7}\\
\hyperref[bg:A.4]{A.4} & functor; functorial invariant & Categories & \thmref{Obs.}{I.4}(c), \thmref{Obs.}{I.6}\\
\hyperref[bg:A.5]{A.5} & natural transformation; naturality & Categories & \thmref{Thm}{III.3}, \thmref{Prop.}{V.4}\\
\hyperref[bg:A.6]{A.6} & slice category & Categories & Table~\ref{tab:notation}, \thmref{Def.}{I.2}\\
\hyperref[bg:A.7]{A.7} & coslice; pointed objects; opposite category & Categories & \thmref{Def.}{I.2}, \thmref{Rem.}{I.3}\\
\hyperref[bg:A.8]{A.8} & the pointed slice $\Meth(\Theta,\theta_0)$ & Categories & Table~\ref{tab:notation}, \thmref{Def.}{I.2}\\
\hyperref[bg:A.9]{A.9} & initial and terminal objects; universal property & Limits & \thmref{Obs.}{I.4}(a), Figure~\ref{fig:interval}\\
\hyperref[bg:A.10]{A.10} & product & Limits & \thmref{Obs.}{I.4}(b), \thmref{Prop.}{I.8}\\
\hyperref[bg:A.11]{A.11} & pullback; fibre product & Limits & \thmref{Obs.}{I.4}(b), \thmref{Prop.}{I.8}\\
\hyperref[bg:A.12]{A.12} & limit; creation of limits & Limits & \thmref{Obs.}{I.4}(b)\\
\hyperref[bg:A.13]{A.13} & adjunction & Limits & \thmref{Obs.}{I.4}(b)\\
\hyperref[bg:A.14]{A.14} & two standard facts about limits & Limits & \thmref{Obs.}{I.4}(b), \S\ref{sec:slice-interval}\\
\hyperref[bg:A.15]{A.15} & monomorphism; sub-object & Limits & \thmref{Def.}{I.2}, \thmref{Cor.}{III.13}\\
\hyperref[bg:A.16]{A.16} & posets; meets; the expressivity order & Limits & \thmref{Obs.}{I.4}(c), \thmref{Prop.}{I.8}(b)\\
\hyperref[bg:A.17]{A.17} & automorphism of a category; transport of structure & Base change & \thmref{Prop.}{II.13}, Exposé~\ref{sec:cut}\\
\hyperref[bg:A.18]{A.18} & lift; lifting problem; obstruction & Base change & \thmref{Obs.}{I.4}(c), \thmref{Obs.}{I.6}\\
\hyperref[bg:A.19]{A.19} & base change; pushforward $\phi_!$; the relative point of view & Base change & Exposés~\ref{sec:thesis}--\ref{sec:slice}, \thmref{Def.}{V.1}\\
\hyperref[bg:A.20]{A.20} & constructions the paper names and leaves unused & Base change & Exposé~\ref{sec:thesis} (where the category theory earns its keep), Exposé~\ref{sec:analogy} (seven borrowed words)\\
\hyperref[bg:B.1]{B.1} & monoidal and symmetric monoidal categories & Monoidal & \thmref{Obs.}{I.7}, \thmref{Thm}{II.10}\\
\hyperref[bg:B.2]{B.2} & direct sum and tensor product in $\FinVect$ & Monoidal & \thmref{Thm}{II.10}, Exposé~\ref{sec:cut}\\
\hyperref[bg:B.3]{B.3} & string diagrams & Monoidal & Exposé~\ref{sec:cut}, \thmref{Thm}{II.10}\\
\hyperref[bg:B.4]{B.4} & cartesian categories; copy and discard & Monoidal & \thmref{Def.}{I.1}, \thmref{Thm}{VI.3}\\
\hyperref[bg:B.5]{B.5} & compact closed categories; cups, caps and wire-bending & Monoidal & \thmref{Thm}{II.11}, \thmref{Thm}{VI.3}\\
\hyperref[bg:B.6]{B.6} & Frobenius algebras and spiders & Monoidal & Exposé~\ref{sec:cut}, \thmref{Thm}{II.10}\\
\hyperref[bg:B.7]{B.7} & tensor networks, contraction, bond dimension & Tensor networks & \thmref{Thm}{II.10}, Exposé~\ref{sec:cut}\\
\hyperref[bg:B.8]{B.8} & the min-cut bound on rank & Tensor networks & \thmref{Thm}{II.10}, Exposé~\ref{sec:cut}\\
\hyperref[bg:B.9]{B.9} & tensor trains, Tucker and MPO; gauge on bonds & Tensor networks & \thmref{Prop.}{II.12}, Exposé~\ref{sec:cut}\\
\hyperref[bg:B.10]{B.10} & Kronecker product, vec, tensor product of subspaces & Tensor networks & \thmref{Thm}{II.11}, Exposé~\ref{sec:cut}\\
\hyperref[bg:B.11]{B.11} & Van Loan--Pitsianis rearrangement; Kronecker rank & Tensor networks & \thmref{Thm}{II.11}, \thmref{Prop.}{II.13}\\
\hyperref[bg:B.12]{B.12} & Schmidt and operator-Schmidt decompositions & Tensor networks & \thmref{Thm}{II.11}, Figure~\ref{fig:cut}\\
\hyperref[bg:B.13]{B.13} & rank of a Hadamard product; face-splitting product & Tensor networks & \thmref{Prop.}{II.12}, \thmref{Thm}{II.10}\\
\hyperref[bg:B.14]{B.14} & sums of methods; Minkowski sums, differences and translates & Para, lenses & \thmref{Obs.}{I.7}, \thmref{Thm}{II.1}\\
\hyperref[bg:B.15]{B.15} & strong monoidal functors; $(\N,+)$ and the permutation groupoid & Para, lenses & \thmref{Obs.}{I.7}, Exposé~\ref{sec:analogy} (the referee record)\\
\hyperref[bg:B.16]{B.16} & bicategories and the Para construction & Para, lenses & \thmref{Def.}{I.1}, \thmref{Rem.}{I.3}\\
\hyperref[bg:B.17]{B.17} & lenses, reverse derivatives and $\Para(\mathsf R)$ & Para, lenses & Exposé~\ref{sec:lens}, \thmref{Def.}{IV.1}\\
\hyperref[bg:C.1]{C.1} & smooth manifolds and smooth maps & Manifolds & \thmref{Def.}{I.2}, \thmref{Def.}{I.5}\\
\hyperref[bg:C.2]{C.2} & tangent spaces, differentials and covectors & Manifolds & Table~\ref{tab:notation}, \thmref{Def.}{I.5}\\
\hyperref[bg:C.3]{C.3} & rank, expressive dimension and first-order image & Manifolds & Table~\ref{tab:notation}, \thmref{Def.}{I.5}\\
\hyperref[bg:C.4]{C.4} & immersions, submersions and embedded submanifolds & Manifolds & \thmref{Obs.}{I.4}, \thmref{Prop.}{IV.3}\\
\hyperref[bg:C.5]{C.5} & constant-rank theorem & Manifolds & \thmref{Def.}{I.5}, \thmref{Prop.}{II.12}\\
\hyperref[bg:C.6]{C.6} & fibres and fibre dimension & Manifolds & Exposé~\ref{sec:fibres}, \thmref{Def.}{I.5}\\
\hyperref[bg:C.7]{C.7} & fibre products, transversality and clean intersection & Manifolds & \thmref{Obs.}{I.4}, \thmref{Prop.}{I.8}\\
\hyperref[bg:C.8]{C.8} & germs & Singularities & \thmref{Thm}{II.7}, Exposé~\ref{sec:image} (the image germ)\\
\hyperref[bg:C.9]{C.9} & smooth and singular points & Singularities & \thmref{Def.}{I.5}, \thmref{Thm}{II.2}\\
\hyperref[bg:C.10]{C.10} & cones and vertices & Singularities & \thmref{Def.}{I.5}, \thmref{Thm}{II.1}\\
\hyperref[bg:C.11]{C.11} & first-order image versus image & Singularities & \thmref{Def.}{I.5}, \thmref{Obs.}{I.6}\\
\hyperref[bg:C.12]{C.12} & the determinantal variety & Varieties & Table~\ref{tab:notation}, \thmref{Thm}{II.1}\\
\hyperref[bg:C.13]{C.13} & the skew-determinantal cone & Varieties & \thmref{Thm}{II.7}, \S\ref{sec:erlangen-hra}\\
\hyperref[bg:C.14]{C.14} & Segre and toric cones & Varieties & \thmref{Prop.}{II.8}, \thmref{Prop.}{II.12}\\
\hyperref[bg:C.15]{C.15} & the Grassmannian & Varieties & \thmref{Thm}{II.2}, \thmref{Thm}{II.7}\\
\hyperref[bg:C.16]{C.16} & images of smooth maps need not be manifolds & Tame geometry & \thmref{Obs.}{I.4}, \thmref{Thm}{II.2}\\
\hyperref[bg:C.17]{C.17} & semialgebraic and definable sets & Tame geometry & Exposé~\ref{sec:thesis} (where the category theory earns its keep), \S\ref{sec:atlas-rules} (arrow rules)\\
\hyperref[bg:C.18]{C.18} & Tarski--Seidenberg and definable dimension & Tame geometry & \thmref{Def.}{I.5}, \S\ref{sec:atlas-rules} (arrow rules)\\
\hyperref[bg:C.19]{C.19} & stratification and local dimension & Tame geometry & \thmref{Thm}{II.2}, \thmref{Thm}{II.7}\\
\hyperref[bg:C.20]{C.20} & moduli sets and moduli spaces & Tame geometry & \thmref{Thm}{II.2}, Exposé~\ref{sec:rosetta}\\
\hyperref[bg:C.21]{C.21} & generic properties & Tame geometry & \thmref{Thm}{II.2}, \thmref{Prop.}{II.8}\\
\hyperref[bg:C.22]{C.22} & real-analytic functions and their zero sets & Tame geometry & \thmref{Lem.}{VI.2}, \thmref{Thm}{VI.5}\\
\hyperref[bg:D.1]{D.1} & Lie groups; the matrix groups of this paper & Lie theory & Table~\ref{tab:notation}, \thmref{Def.}{I.5}\\
\hyperref[bg:D.2]{D.2} & tori, diagonal monoids and saturation & Lie theory & \thmref{Thm}{II.5}, \thmref{Prop.}{II.8}\\
\hyperref[bg:D.3]{D.3} & permutations, monomial matrices, semidirect products, normalisers & Lie theory & \thmref{Prop.}{II.8}, \thmref{Prop.}{II.9}\\
\hyperref[bg:D.4]{D.4} & connectedness and the identity component & Lie theory & \thmref{Thm}{V.3}, \S\ref{sec:erlangen-invariants}\\
\hyperref[bg:D.5]{D.5} & Lie algebras; the commutator; $\gl$, $\so$, $\Sym$, $\wedge$ & Lie theory & \thmref{Thm}{II.7}, \thmref{Thm}{III.5}\\
\hyperref[bg:D.6]{D.6} & the exponential map; one-parameter subgroups & Lie theory & \thmref{Thm}{III.5}, \thmref{Prop.}{III.7}\\
\hyperref[bg:D.7]{D.7} & the Cayley map and its Neumann truncation & Lie theory & \thmref{Cor.}{II.6}, \thmref{Thm}{II.7}\\
\hyperref[bg:D.8]{D.8} & Lie subgroups and subalgebras; the generated Lie algebra & Lie theory & \thmref{Thm}{V.3}, Figure~\ref{fig:rebase}\\
\hyperref[bg:D.9]{D.9} & Yamabe's theorem & Lie theory & \thmref{Thm}{V.3}, \S\ref{sec:base-rebase}\\
\hyperref[bg:D.10]{D.10} & monoids, units, monoid orbits, homomorphisms & Lie theory & \thmref{Thm}{V.3}, \thmref{Prop.}{II.9}\\
\hyperref[bg:D.11]{D.11} & Cartan--Dieudonné theorem; Scherk's sharp form & Lie theory & \thmref{Thm}{II.7}, \S\ref{sec:erlangen-hra}\\
\hyperref[bg:D.12]{D.12} & group actions & Invariants & \thmref{Def.}{I.2}, Table~\ref{tab:notation}\\
\hyperref[bg:D.13]{D.13} & orbits, stabilisers, free actions & Invariants & \thmref{Thm}{II.5}, \thmref{Prop.}{II.9}\\
\hyperref[bg:D.14]{D.14} & Haar measure; invariant measures on homogeneous spaces & Invariants & \thmref{Thm}{II.2}, \S\ref{sec:strata-three}\\
\hyperref[bg:D.15]{D.15} & invariants and complete invariants & Invariants & \thmref{Thm}{II.5}, \S\ref{sec:erlangen-invariants}\\
\hyperref[bg:D.16]{D.16} & first fundamental theorem for $\Ort(n)$; Witt's extension theorem & Invariants & \thmref{Thm}{II.5}, Exposé~\ref{sec:related}\\
\hyperref[bg:D.17]{D.17} & Klein's Erlangen programme & Invariants & \thmref{Def.}{II.14}, Exposé~\ref{sec:erlangen}\\
\hyperref[bg:D.18]{D.18} & equivariance and invariance of maps, update rules and merge rules & Invariants & \thmref{Def.}{II.14}, \thmref{Obs.}{II.15}\\
\hyperref[bg:D.19]{D.19} & multihomogeneous maps; multidegree & Invariants & \thmref{Thm}{III.10}, Exposé~\ref{sec:image} (the image germ)\\
\hyperref[bg:D.20]{D.20} & orbit space of a positive torus & Invariants & \thmref{Thm}{III.10}, \thmref{Cor.}{III.11}\\
\hyperref[bg:D.21]{D.21} & gauge symmetry of a parametrisation & Gauge, metrics & \thmref{Def.}{I.5}, Table~\ref{tab:notation}\\
\hyperref[bg:D.22]{D.22} & quotients and descent; the quotient manifold theorem & Gauge, metrics & \thmref{Prop.}{V.4}, \thmref{Prop.}{III.4}\\
\hyperref[bg:D.23]{D.23} & metrics, cometrics and the pulled-back cometric & Gauge, metrics & \thmref{Def.}{III.1}, \thmref{Obs.}{III.2}\\
\hyperref[bg:D.24]{D.24} & Riemannian gradient flow; tangent projection; retraction & Gauge, metrics & \thmref{Prop.}{III.4}, \thmref{Thm}{III.12}\\
\hyperref[bg:D.25]{D.25} & distributions, Lie brackets and Frobenius' theorem & Gauge, metrics & \thmref{Prop.}{IV.3}, \S\ref{sec:lens-integrability}\\
\hyperref[bg:D.26]{D.26} & Noether's theorem for gradient flows & Conservation & \thmref{Thm}{III.5}, Figure~\ref{fig:noether}\\
\hyperref[bg:D.27]{D.27} & Cartan decomposition; polar decomposition & Conservation & \thmref{Thm}{III.5}, \thmref{Prop.}{III.7}\\
\hyperref[bg:D.28]{D.28} & balancedness; the conserved $B\T B-AA\T$ & Conservation & \thmref{Prop.}{III.7}, \thmref{Prop.}{III.8}\\
\hyperref[bg:D.29]{D.29} & the moment map & Conservation & \thmref{Prop.}{III.7}, \thmref{Cor.}{III.6}\\
\hyperref[bg:D.30]{D.30} & Kempf--Ness theorem, real form & Conservation & \thmref{Prop.}{III.7}, \thmref{Cor.}{III.6}\\
\hyperref[bg:D.31]{D.31} & the closed orbit in each fibre of $(B,A)\mapsto BA$ & Conservation & \thmref{Prop.}{V.4}, \S\ref{sec:base-combine}\\
\hyperref[bg:D.32]{D.32} & Ky Fan's maximum principle & Matrix tools & \thmref{Prop.}{II.3}, \S\ref{sec:erlangen-loft}\\
\hyperref[bg:D.33]{D.33} & Sylvester equations & Matrix tools & \thmref{Thm}{III.12}, Figure~\ref{fig:gauge}\\
\hyperref[bg:D.34]{D.34} & the perspective map; convex combinations & Matrix tools & \thmref{Thm}{VI.5}, Figure~\ref{fig:prefix}\\
\end{longtable}}

% ======================================================================================================================
\subsection*{Part A. Categories without tears}
\phantomsection\label{app:bg-a}\addcontentsline{toc}{subsection}{Part A. Categories without tears}

Part A collects the category theory the exposés use, from the definition of a category up to base change, and
illustrates every notion in one running example, the pointed slice $\Meth(\Theta,\theta_0)$ of weight-space methods.
Entries \bgref{A.1} to \bgref{A.8} set up categories, functors and slices, \bgref{A.9} to \bgref{A.16} cover universal
properties, limits and orders, and \bgref{A.17} to \bgref{A.19} cover transport, lifting and base change, with
\bgref{A.20} listing the constructions the paper names and leaves unused.

\subsubsection*{Categories and slices}

\begin{bgentry}{A.1}{category}
A \textbf{category} $\mathcal C$ has objects and, for any two objects $X,Y$, a set $\mathcal C(X,Y)$ of
\textbf{morphisms}, or arrows, $f:X\to Y$. Arrows $f:X\to Y$ and $g:Y\to Z$ have a composite $g\circ f:X\to Z$.
Composition is associative, and every object has an identity arrow $\id_X$ with $f\circ\id_X=f=\id_Y\circ f$.
\end{bgentry}

The paper uses $\Set$ (sets and functions), $\Diff$ (smooth manifolds and smooth maps), $\FinVect$ (finite-dimensional
real vector spaces and linear maps, where tensor networks are evaluated) and the pointed categories $\Set_*$ and
$\Diff_*$ of \bgref{A.7}. In $\Meth(\Theta,\theta_0)$ the objects are methods $M=(Q,q_0,\rho)$, and an arrow $h:M\to N$
is a pointed smooth map $h:Q_M\to Q_N$ with $\rho_N\circ h=\rho_M$, read ``$N$ simulates $M$''. Simulations compose,
since $\rho_P\circ k\circ h=\rho_N\circ h=\rho_M$ for $h:M\to N$ and $k:N\to P$. For instance, $h(B)=(B,A_0)$ is an
arrow from LoRA-FA$(A_0)$ \citep{zhang2023lora} to LoRA pointed at $(0,A_0)$. The simulation arrows of the atlas are
arrows of this category, each recorded with its map $h$.

\fbAusedin{Table~\ref{tab:notation}, \thmref{Definition}{I.2}, Exposés~\ref{sec:thesis}--\ref{sec:slice},
\S\ref{sec:atlas-arrows}, Figure~\ref{fig:graph}. \textsf{Reading:} \citet[Ch.~1]{xleinster2014basic}.}

\begin{bgentry}{A.2}{commutative diagram}
A \textbf{diagram} is a family of objects joined by arrows, drawn as a directed graph. It \textbf{commutes} if any two
directed paths with the same start and the same end compose to the same arrow. A triangle commutes when one side equals
the composite of the other two, and a square commutes when its two routes between opposite corners agree.
\end{bgentry}

A simulation is a commuting triangle over $\Theta$, whose two paths from $Q_M$ to $\Theta$ are
\begin{equation*}
  Q_M\xrightarrow{\;h\;}Q_N\xrightarrow{\;\rho_N\;}\Theta\qquad\text{and}\qquad Q_M\xrightarrow{\;\rho_M\;}\Theta.
\end{equation*}
For $M=\mathrm{FA}(A_0)$, $N$ LoRA pointed at $(0,A_0)$ and $h(B)=(B,A_0)$, both paths send $B$ to the weight
$W_0+BA_0$ (scale absorbed into $B$), so the triangle commutes. Figure~\ref{fig:interval}(a) draws, for a general
method, the triangle of Figure~\ref{fig:bg-interval}, which commutes only up to a pointing defect when the method starts
away from $\theta_0$ (\bgref{A.19}). Commuting squares, used for merging on extensions, return in \bgref{A.18}.

\fbAusedin{Figure~\ref{fig:interval}, \thmref{Definition}{I.2}, \thmref{Definition}{VI.1}, \thmref{Definition}{V.1},
Exposé~\ref{sec:base}. \textsf{Reading:} \citet[\S1.1]{xleinster2014basic}.}

\begin{bgentry}{A.3}{isomorphism}
An arrow $f:X\to Y$ is an \textbf{isomorphism}, written $X\cong Y$, if some $g:Y\to X$ satisfies $g\circ f=\id_X$ and
$f\circ g=\id_Y$. The inverse $g$ is then unique. In $\Diff$ the isomorphisms are the diffeomorphisms, so a smooth
bijection need not be one, as $t\mapsto t^3$ shows. In $\Meth$ an isomorphism $M\cong N$ is a pointed diffeomorphism
$h:Q_M\to Q_N$ with $\rho_N\circ h=\rho_M$. Its inverse lies over $\Theta$ automatically, since
$\rho_M\circ h^{-1}=\rho_N$.
\end{bgentry}

The LoRA gauge (\bgref{D.21}) $(B,A)\mapsto(Bg^{-1},gA)$, $g\in\GL_r$, is an isomorphism from LoRA pointed at
$(B_0,A_0)$ to LoRA pointed at $(B_0g^{-1},gA_0)$. Other examples are rsLoRA \citep{kalajdzievski2023rank} $\cong$ LoRA
by $(B,A)\mapsto(\sqrt r\,B,A)$, and $\mathrm{LoRA}_r\oplus\mathrm{LoRA}_s\cong\mathrm{LoRA}_{r+s}$ by concatenating
factors (\bgref{B.14}). Isomorphic methods have the same image, the same $S_1$ and the same first-order deficit
(\bgref{C.3}, \thmref{Observation}{I.6}). They can still train differently, because the optimizer acts on the
coordinates of $Q$. For LoRA's gauge with block-scalar rates (one rate per factor), gradient flow is unchanged exactly
when $g\in\Ort(r)$, at points where $A$ has full row rank and $r<n$ (\thmref{Theorem}{III.3}; \bgref{D.23}).

\fbAusedin{\thmref{Observation}{I.6}, \thmref{Observation}{I.7}, \S\ref{sec:strata-three}, \thmref{Theorem}{II.2},
\thmref{Proposition}{VI.4}, \S\ref{sec:atlas-arrows}.}

\begin{bgentry}{A.4}{functor; functorial invariant}
A \textbf{functor} $F:\mathcal C\to\mathcal D$ sends each object $X$ to an object $FX$ and each arrow $f:X\to Y$ to an
arrow $Ff:FX\to FY$, with $F(g\circ f)=Fg\circ Ff$ and $F(\id_X)=\id_{FX}$. A functor sends isomorphisms to
isomorphisms, so $FX$ is an isomorphism invariant of $X$. A \textbf{forgetful} functor discards structure, as
$\Diff_*\to\Diff$ discards the base point.
\end{bgentry}

The image is a functor $\Im$ from $\Meth(\Theta,\theta_0)$ to the subsets of $\Theta$ containing $\theta_0$, ordered by
inclusion (a poset, \bgref{A.16}), because an arrow $h:M\to N$ gives $\Im M=\rho_N(h(Q_M))\subseteq\Im N$. In the same
way the chain rule at $q_0$ makes $S_1$ a functor to the linear subspaces of $\Theta$, with $S_1(M)\subseteq S_1(N)$
along every pointed arrow. Read backwards, functoriality gives the atlas test O1 and the first half of O3
(\S\ref{sec:atlas-tests}). If $\Im M\not\subseteq\Im N$ or $S_1(M)\not\subseteq S_1(N)$, then $N$ does not simulate
$M$. Base change $\phi_!$ (\bgref{A.19}) and transport $T_*$ (\bgref{A.17}) are functors between slices.

\fbAusedin{\thmref{Observation}{I.4}(c), \thmref{Observation}{I.6}, \S\ref{sec:atlas-arrows},
\thmref{Proposition}{II.13}, \thmref{Definition}{V.1}, \S\ref{sec:lens-functor}. \textsf{Reading:}
\citet[\S1.2]{xleinster2014basic}.}

\begin{bgentry}{A.5}{natural transformation; naturality}
For functors $F,G:\mathcal C\to\mathcal D$, a \textbf{natural transformation} $\alpha:F\Rightarrow G$ is a family of
arrows $\alpha_X:FX\to GX$ with $Gf\circ\alpha_X=\alpha_Y\circ Ff$ for every $f:X\to Y$. This naturality square says that
applying $\alpha$ and then moving along $f$ agrees with moving first and applying $\alpha$ after. A group $K$ can be
viewed as a category with one object whose arrows are the elements of $K$. A functor from it to $\Set$ is then a set
with a $K$-action, and a natural transformation between two such functors is exactly a $K$-equivariant map.
\end{bgentry}

The copy map $\Delta_X(x)=(x,x)$ is natural, $\Delta_Y\circ f=(f\times f)\circ\Delta_X$, which is rule R7 of the merge
calculus (\bgref{B.4}). Elsewhere in the paper, natural for a group of changes of chart means equivariant in this sense
(\bgref{D.18}). Task arithmetic \citep{ilharco2022editing} is natural for all of $\GL(\Theta)$, and DARE
\citep{yu2023language} for invertible diagonal matrices at a fixed mask (\thmref{Proposition}{V.4}(c)). Gradient flow
on LoRA commutes with the gauge $g$ exactly when the induced cometrics (\bgref{D.23}) agree, which for block-scalar rates
singles out $\Ort(r)$ (\thmref{Theorem}{III.3}).

\fbAusedin{Exposé~\ref{sec:thesis} (where the category theory earns its keep), \thmref{Theorem}{III.3},
\S\ref{sec:fibres-naturality}, \thmref{Proposition}{V.4}, \S\ref{sec:merge-calculus}, \thmref{Theorem}{VI.3}.
\textsf{Reading:} \citet[\S1.3]{xleinster2014basic}.}

\begin{bgentry}{A.6}{slice category}
For an object $X$ of $\mathcal C$, the \textbf{slice} $\mathcal C/X$ has as objects the arrows $a:A\to X$, called
objects over $X$, and as morphisms $(A,a)\to(B,b)$ the arrows $h:A\to B$ of $\mathcal C$ with $b\circ h=a$. Composition
and identities are those of $\mathcal C$. The object $X$ is the \textbf{base}, and $a$ is the \textbf{structure map}.
\end{bgentry}

$\Diff/\Theta$ consists of smooth maps $\rho:Q\to\Theta$ into weight space without base points, and its morphisms are
the weak morphisms of \thmref{Definition}{I.2}, which need not send $q_0$ to $q_0'$. A slice always has a terminal
object, the identity $(X,\id_X)$, which here is full fine-tuning. Working in a slice is the relative point of view of
\bgref{A.19}, in which an object is studied together with its map to a base.

\fbAusedin{Table~\ref{tab:notation}, \thmref{Definition}{I.2}, Exposé~\ref{sec:thesis} (where the category theory earns
its keep), Rosetta entries 3 and 4 (\S\ref{sec:rosetta-six}), the seven borrowed words of Exposé~\ref{sec:analogy}.
\textsf{Reading:} \citet[\S II.6]{xmaclane1998categories}.}

\begin{bgentry}{A.7}{coslice; pointed objects; opposite category}
The \textbf{coslice} $X/\mathcal C$ has as objects the arrows $x:X\to A$ out of $X$, and as morphisms $(A,x)\to(B,x')$
the arrows $h:A\to B$ with $h\circ x=x'$. For the one-point space $X=\mathrm{pt}$, an arrow $\mathrm{pt}\to Q$ is a point
$q_0$, so $\Diff_*:=\mathrm{pt}/\Diff$ is the category of \textbf{pointed manifolds} $(Q,q_0)$ and \textbf{pointed
maps}, those with $h(q_0)=q_0'$. Likewise $\Set_*:=\mathrm{pt}/\Set$. The \textbf{opposite category}
$\mathcal C^{\mathrm{op}}$ has the same objects and every arrow reversed, and
$(\mathcal C/X)^{\mathrm{op}}\cong X/\mathcal C^{\mathrm{op}}$.
\end{bgentry}

A method sends its initialisation to the checkpoint, $\rho(q_0)=\theta_0$, so it is an arrow of $\Diff_*$, and a weak
morphism forgets the base points and lives in $\Diff$. Opposites enter in \thmref{Remark}{I.3}. A simulation
$h:M\to N$ runs from $Q_M$ to $Q_N$. The network with $Q_M$-parameters is obtained from the one with $Q_N$-parameters by
precomposing with $h$,
\begin{equation*}
  f\circ(\rho_M\times\id_X)=f\circ(\rho_N\times\id_X)\circ(h\times\id_X),
\end{equation*}
with $X$ the input space. Read as a change of parameters for networks, $h$ therefore points the other way, from the
$Q_N$-network to the $Q_M$-network, and \thmref{Remark}{I.3} records this reversal with an opposite category
(\bgref{B.16}).

\fbAusedin{\thmref{Definition}{I.2}, \thmref{Remark}{I.3}, \thmref{Observation}{I.4}, \S\ref{sec:slice-para}.}

\begin{bgentry}{A.8}{the pointed slice $\Meth(\Theta,\theta_0)$}
Combining \bgref{A.6} and \bgref{A.7}, $\Meth(\Theta,\theta_0):=\Diff_*/(\Theta,\theta_0)$. An object is a pointed smooth
map $\rho:(Q,q_0)\to(\Theta,\theta_0)$, and a morphism $h:M\to N$ is a pointed smooth map with $\rho_N\circ h=\rho_M$.
Equivalently, an object is a factorisation $\mathrm{pt}\xrightarrow{\,q_0\,}Q\xrightarrow{\,\rho\,}\Theta$ of the point
$\theta_0:\mathrm{pt}\to\Theta$, and a morphism is a map between the middle terms that is compatible with both ends.
\end{bgentry}

On one box, LoRA is $\big(\R^{m\times r}\times\R^{r\times n},\,(0,A_0),\,(B,A)\mapsto W_0+sBA\big)$, and OFT
\citep{qiu2023controlling} is $Q\mapsto W_0\,C(Q)$ on skew matrices (block-diagonal in OFT), pointed at $0$, with $C$
the Cayley map (\bgref{D.7}). Full fine-tuning is $(\Theta,\theta_0,\id)$ and the frozen model is
$(\mathrm{pt},*,\theta_0)$. The network $f$ appears nowhere in the definition, so the category depends only on
$(\Theta,\theta_0)$ (\thmref{Remark}{I.3}). The factorisation picture already shows why the frozen model and full
fine-tuning sit at the two ends (\bgref{A.9}).

\fbAusedin{Table~\ref{tab:notation}, \thmref{Definition}{I.2}, \S\ref{sec:slice-methods}, \thmref{Remark}{I.3}, Rosetta
entries 3 and 4 (\S\ref{sec:rosetta-six}).}

\subsubsection*{Universal properties}

\begin{bgentry}{A.9}{initial and terminal objects; universal property}
An object $I$ is \textbf{initial} if every object $X$ receives exactly one arrow $I\to X$, and $T$ is \textbf{terminal}
if every $X$ has exactly one arrow $X\to T$. Each is unique up to a unique isomorphism when it exists. Characterising an
object by the arrows into or out of it in this way is called a \textbf{universal property}. In $\Set$ the empty set is
initial and a one-point set is terminal, and in $\Set_*$ and $\Diff_*$ the one-point space is both.
\end{bgentry}

In $\Meth(\Theta,\theta_0)$ the frozen model $(\mathrm{pt},*,\theta_0)$ is initial. The only pointed map
$\mathrm{pt}\to Q$ is $*\mapsto q_0$, and it lies over $\Theta$ because $\rho(q_0)=\theta_0$. Full fine-tuning
$(\Theta,\theta_0,\id)$ is terminal, since an arrow $h$ from $M$ must satisfy $\id\circ h=\rho_M$. So the unique arrow
from $M$ to Full FT is $\rho_M$ itself, the merge, and every weight-space method that starts at $\theta_0$ sits in the
commuting triangle of Figure~\ref{fig:bg-interval}.

\begin{figure}[tbp]
\centering
\begin{tikzpicture}[>={Stealth[length=2.3mm,width=1.8mm]},
  obj/.style={align=center, inner sep=3pt, text=ink},
  arr/.style={->, draw=ink, line width=0.6pt, shorten >=1pt, shorten <=1pt},
  lab/.style={font=\small, text=ink}]
  \node[obj] (M) at (0,1.95) {method $M$\\[1pt] $(Q,q_0)$};
  \node[obj] (F) at (-3.3,0) {Frozen\\[1pt] $(\mathrm{pt},*)$};
  \node[obj] (T) at (3.3,0) {Full FT\\[1pt] $(\Theta,\theta_0)$};
  \draw[arr] (F) -- node[above left=-1pt, lab] {$q_0$} (M);
  \draw[arr] (M) -- node[above right=-1pt, lab] {$\rho$, the merge} (T);
  \draw[arr] (F) -- node[below, lab] {$\theta_0$} (T);
\end{tikzpicture}
\caption{The interval of $\Meth(\Theta,\theta_0)$. The left arrow is the only arrow out of the initial object, the right
  arrow $\rho$ is the only arrow into the terminal object, and the triangle commutes because $\rho(q_0)=\theta_0$.
  Figure~\ref{fig:interval}(a) draws the same triangle in the main text, where a pointing defect
  $\rho_M(q_0)\neq\theta_0$ leaves it commuting only up to that defect.}
\label{fig:bg-interval}
\end{figure}

\fbAusedin{\thmref{Observation}{I.4}(a), \S\ref{sec:slice-interval}, Figure~\ref{fig:interval}, Rosetta entries 3 and 4
(\S\ref{sec:rosetta-six}), rule A1 in \S\ref{sec:atlas-rules}, \S\ref{sec:base-next}.}

\begin{bgentry}{A.10}{product}
A \textbf{product} of $X$ and $Y$ is an object $X\times Y$ with projections $p_X:X\times Y\to X$ and
$p_Y:X\times Y\to Y$ such that for every object $T$ and arrows $u:T\to X$, $v:T\to Y$ there is exactly one arrow
$\langle u,v\rangle:T\to X\times Y$ with $p_X\circ\langle u,v\rangle=u$ and $p_Y\circ\langle u,v\rangle=v$. It is unique
up to unique isomorphism, and it need not exist. In a slice $\mathcal C/Z$, the product of $(X,a)$ and $(Y,b)$ is the
pullback of $a$ and $b$ in $\mathcal C$ (\bgref{A.11}).
\end{bgentry}

In $\Meth$ a product $M\times N$ has arrows to $M$ and to $N$, so both factors simulate it, and every method $T$ with
arrows $u:T\to M$ and $v:T\to N$ maps to it by exactly one arrow compatible with $u$ and $v$. When it exists, its image is
$\Im M\cap\Im N$. For the one-sided LoRAs with frozen frames $B_0$ of full column rank and $A_0$ of full row rank,
\begin{equation*}
  \mathrm{FA}(A_0)\times\mathrm{FB}(B_0)\;\cong\;\big(\R^{r\times r},\,0,\,R\mapsto W_0+B_0RA_0\big),
\end{equation*}
with projections $R\mapsto B_0R$ and $R\mapsto RA_0$. In the SVD frames of $W_0$, when $\rk W_0\ge r$, this is LoRA-XS
\citep{baazy2024lora} (\thmref{Proposition}{I.8}(a)). The sum $\oplus$ of \thmref{Observation}{I.7} is a different
operation, whose image is a Minkowski sum (\bgref{B.14}).

\fbAusedin{\thmref{Observation}{I.4}(b), \thmref{Proposition}{I.8}, \S\ref{sec:slice-products},
\thmref{Proposition}{II.3}, rule A4 in \S\ref{sec:atlas-rules}, \S\ref{sec:rosetta-reverse}. \textsf{Reading:}
\citet[\S5.1]{xleinster2014basic}.}

\begin{bgentry}{A.11}{pullback; fibre product}
Given $f:X\to Z$ and $g:Y\to Z$, a \textbf{pullback} is an object $P$ with arrows $p:P\to X$ and $q:P\to Y$ such that
$f\circ p=g\circ q$, and such that whenever $u:T\to X$ and $v:T\to Y$ satisfy $f\circ u=g\circ v$, exactly one
$t:T\to P$ has $p\circ t=u$ and $q\circ t=v$ (Figure~\ref{fig:bg-pullback}). In $\Set$ the pullback is the
\textbf{fibre product} $X\times_ZY=\{(x,y):f(x)=g(y)\}$. In $\Diff$, if this set is an embedded submanifold of
$X\times Y$, it is the pullback, and transversality (\bgref{C.7}) is the standard sufficient condition for that. A
pullback in $\Diff$ can fail to exist.
\end{bgentry}

Products in $\Meth$ are fibre products $Q_M\times_\Theta Q_N=\{(q,q'):\rho_M(q)=\rho_N(q')\}$. For $\mathrm{FA}(A_0)$
and $\mathrm{FB}(B_0)$, with the ranks of \bgref{A.10}, this set is the linear space $\{(B_0R,RA_0)\}$, so the product
exists. For $\mathrm{LoRA}_1$ and the frozen model on $1\times1$ weights the set is the cross $\{ba=0\}\subset\R^2$, no
manifold structure on it has the universal property, and the product does not exist (\bgref{A.13}). Pulling back along
the point $\theta_0:\mathrm{pt}\to\Theta$ gives the fibre $\rho^{-1}(\theta_0)$, the underlying set of
$M\times\text{Frozen}$ whenever that product exists.

\paragraph{HRA is a meet of images, not a product}
Take $W_0$ injective and $r$ even with $2\le r\le n-2$. The fibre product of $\mathrm{LoRA}_r$ with the rotation orbit
$W_0\cdot\SO(n)$ has local dimension $d(\mathrm{HRA}_r)+r^2=rn+r(r-1)/2$ over generic points of the meet
(\bgref{C.19}). This exceeds $rn=\lvert\mathrm{HRA}_r\rvert$ because $r\ge2$, so HRA \citep{yuan2024bridging} is not the
product of the two methods, and \thmref{Theorem}{II.7} keeps only the meet of their images (\bgref{A.16}).

\begin{figure}[tbp]
\centering
\begin{tikzcd}[column sep=3.2em, row sep=2.8em, arrows={line width=0.6pt}]
  T \arrow[drr, bend left=22, "v"] \arrow[ddr, bend right=22, "u"'] \arrow[dr, dashed, "\exists!\,t" description]
    & & \\
  & Q_M\times_\Theta Q_N \arrow[r, "p_N"] \arrow[d, "p_M"'] \arrow[dr, phantom, "\lrcorner", very near start]
    & Q_N \arrow[d, "\rho_N"] \\
  & Q_M \arrow[r, "\rho_M"'] & \Theta
\end{tikzcd}
\caption{A pullback square. It commutes, $\rho_M\circ p_M=\rho_N\circ p_N$, and any test object $T$ whose maps $u,v$
  agree over $\Theta$ factors through the corner by exactly one $t$. The small corner mark is the usual sign of a
  pullback.}
\label{fig:bg-pullback}
\end{figure}

\fbAusedin{\thmref{Observation}{I.4}(b), \thmref{Proposition}{I.8}, \thmref{Proposition}{II.3},
\thmref{Theorem}{II.7}(d), \thmref{Corollary}{III.13}.}

\begin{bgentry}{A.12}{limit; creation of limits}
A \textbf{cone} over a diagram $D$ in $\mathcal C$ is an object $L$ with an arrow to every object of $D$, such that each
triangle formed with an arrow of $D$ commutes. A \textbf{limit} of $D$ is a cone through which every other cone factors
by exactly one arrow. This categorical cone is unrelated to the geometric cones of \bgref{C.10}. Terminal objects,
products and pullbacks are the limits of the empty diagram, of two objects, and of two arrows with a common target. A
functor $U:\mathcal C\to\mathcal D$ \textbf{creates limits} if, for every diagram $D$ in $\mathcal C$ and every limit cone
of $U\circ D$, exactly one cone over $D$ is sent by $U$ to that limit cone, and this cone is a limit of $D$.
\end{bgentry}

The forgetful functor $\Diff_*\to\Diff$ creates limits (\bgref{A.14}). Concretely, when the underlying manifolds have a
pullback, the pullback of pointed manifolds is that pullback, pointed at the pair of base points. This is how
\thmref{Observation}{I.4}(b) computes products of methods as fibre products in $\Diff$.

\fbAusedin{\thmref{Observation}{I.4}(b). \textsf{Reading:} \citet[Ch.~V]{xmaclane1998categories};
\citet[Ch.~3]{xriehl2016category}.}

\begin{bgentry}{A.13}{adjunction}
A functor $F:\mathcal C\to\mathcal D$ is \textbf{left adjoint} to $G:\mathcal D\to\mathcal C$, and $G$ is \textbf{right
adjoint} to $F$, written $F\dashv G$, if there are bijections $\mathcal D(FX,Y)\cong\mathcal C(X,GY)$ natural in $X$ and
$Y$. Many forgetful functors are right adjoints, whose left adjoints add the forgotten structure freely.
\end{bgentry}

The forgetful functor $\Diff_*\to\Diff$ is right adjoint to $X\mapsto(X\sqcup\mathrm{pt},\mathrm{pt})$, which adds a
disjoint base point, since a pointed map out of $X\sqcup\mathrm{pt}$ is the same as an arbitrary map out of $X$ (a
manifold may have components of different dimensions). Right adjoints preserve limits (\bgref{A.14}), so a product of
methods, if it exists, is also a pullback in $\Diff$ and can be tested with unpointed maps.
\thmref{Observation}{I.4}(b) uses this to rule out $\mathrm{LoRA}_1\times\text{Frozen}$ on $1\times1$ weights. The axis
curves $t\mapsto(t,0)$ and $t\mapsto(0,t)$ in the cross would lift smoothly through the point over the origin, so the
differential of the projection from the product to the parameter plane $\R^2$ of $\mathrm{LoRA}_1$ would have rank $2$
there. That projection would then be a submersion near the point (\bgref{C.4}), and its image would contain an open set,
which the cross does not.

\fbAusedin{\thmref{Observation}{I.4}(b). \textsf{Reading:} \citet[Ch.~2 and Ch.~6]{xleinster2014basic}.}

\begin{bgfact}{A.14}{two standard facts about limits}
Let $\mathcal C$ be a category.
\begin{enumerate}[label=(\alph*)]
  \item For every object $X$ of $\mathcal C$, the forgetful functor $X/\mathcal C\to\mathcal C$ creates limits. In
    particular $\Diff_*\to\Diff$ does.
  \item A right adjoint preserves every limit that exists in its domain.
\end{enumerate}
\end{bgfact}

\paragraph{The idea behind (a) and (b)}
(a) The structure arrows $X\to D_j$ of a diagram in $X/\mathcal C$ form a cone in $\mathcal C$, which factors through
the limit by exactly one arrow, and that arrow is the structure map of the limit. (b) For $F\dashv G$,
\begin{equation*}
  \mathcal C(X,G\lim D)\cong\mathcal D(FX,\lim D)\cong\lim\mathcal D(FX,D_j)\cong\lim\mathcal C(X,GD_j),
\end{equation*}
naturally in $X$.

\medskip\noindent
\thmref{Observation}{I.4}(b) combines them. Products in $\Meth$ are pullbacks in $\Diff_*$, which by (a) exist whenever
the underlying pullback exists in $\Diff$, and by (b) and \bgref{A.13} only then. The transversality hypothesis of
\bgref{C.7} fails at the base point whenever $\lvert M\rvert+\lvert N\rvert<\dim\Theta$, with $\lvert M\rvert=\dim Q_M$
(\bgref{C.3}), so the products that exist, such as LoRA-XS, come from other structure (\bgref{A.11}).

\fbAusedin{\thmref{Observation}{I.4}(b), \S\ref{sec:slice-interval}. \textsf{Reading:}
\citet[Ch.~V]{xmaclane1998categories}; \citet[Ch.~3 and Ch.~4]{xriehl2016category}.}

\begin{bgentry}{A.15}{monomorphism; sub-object}
An arrow $m:A\to B$ is a \textbf{monomorphism} if $m\circ u=m\circ v$ implies $u=v$ for all arrows $u,v:T\to A$. A
\textbf{sub-object} of $B$ is a monomorphism into $B$, two of them being identified when they differ by an isomorphism
over $B$. Testing with maps out of one- and two-point manifolds shows that the monomorphisms of $\Diff$, and those of
$\Meth$, are exactly the arrows whose underlying map is injective. Such a map need not be an embedding, as
$t\mapsto t^3$ shows.
\end{bgentry}

Freezing coordinates at their base values gives an injective affine map, hence a sub-object, written $\hookrightarrow$ in
the atlas (rule A2). Examples are $\mathrm{FA}(A_0)\hookrightarrow\mathrm{LoRA}$ pointed at $(0,A_0)$, by
$B\mapsto(B,A_0)$, and MiCA \citep{rudiger2026mica} ${}\hookrightarrow{}$ LoRA pointed at $(B_0,0)$, by
$A\mapsto(B_0,A)$. These maps are embeddings, a stronger property than injectivity. The paper writes sub-object and never
monomorphism. In \thmref{Corollary}{III.13}, LoRA-FA is a sub-object of LoRA and so a different object, while random-$A$
LoRA and EVA \citep{paischer2024parameter} are the same map $\rho$ with two different base points in its fibre over
$\theta_0$.

\fbAusedin{\thmref{Definition}{I.2}, rule A2 in \S\ref{sec:atlas-rules}, \thmref{Corollary}{III.13}.}

\begin{bgentry}{A.16}{posets; meets; the expressivity order}
A \textbf{preorder} is a category with at most one arrow between any two objects; write $x\le y$ when there is one. It
is a \textbf{poset} if $x\le y\le x$ forces $x=y$. In a poset the product of $x$ and $y$ is their \textbf{meet}
$x\wedge y$, the greatest lower bound, and for subsets ordered by inclusion the meet is the intersection. Two elements are
\textbf{incomparable} if neither $x\le y$ nor $y\le x$, and $x$ is \textbf{maximal} if $x\le y$ implies $x=y$. Every
category has a \textbf{poset reflection}, in which $X\le Y$ when some arrow $X\to Y$ exists and $X,Y$ are identified when
arrows run both ways. A \textbf{Hasse diagram} draws a finite poset with $y$ above $x$ when $x<y$, joining them by an
edge when nothing lies strictly between.
\end{bgentry}

The poset reflection of $\Meth$, which we call the expressivity order, orders methods by simulation.
Figure~\ref{fig:graph} draws it for the core methods as a layered Hasse diagram, from the frozen model at the bottom to
full fine-tuning at the top. A meet of images always exists and need not be the image of a product. For injective $W_0$
and even $r\le n-2$, the image of HRA without Gram--Schmidt is the meet of LoRA's cone with the rotation orbit
$W_0\cdot\SO(n)$ (\bgref{D.11}), and HRA is not their product (\thmref{Theorem}{II.7}). Split and zero-init LoRA have
incomparable images (\thmref{Theorem}{II.2}(d)), so neither simulates the other. The frontier problem asks for the
maximal methods of budget at most $k$, which are automatically pairwise incomparable.

\fbAusedin{\thmref{Observation}{I.4}(c), \thmref{Proposition}{I.8}(b), \thmref{Theorem}{II.2}, \thmref{Theorem}{II.7},
\S\ref{sec:erlangen-hra}, Figure~\ref{fig:graph}, the open problems of Exposé~\ref{sec:analogy}. \textsf{Reading:}
\citet[Ch.~1 and Ch.~2]{xdavey2002introduction}.}

\subsubsection*{Base change and lifting}

\begin{bgentry}{A.17}{automorphism of a category; transport of structure}
An \textbf{automorphism} of a category $\mathcal C$ is a functor $F:\mathcal C\to\mathcal C$ with an inverse functor, so
that both composites are identity functors. It carries initial objects, products, isomorphisms and every other notion
defined from arrows alone to notions of the same kind. \textbf{Transport of structure} builds such a functor from an
invertible map of the underlying data.
\end{bgentry}

For $T\in\GL(\Theta)$ put $\tau_T(\theta)=\theta_0+T(\theta-\theta_0)$ and $T_*M:=(Q,q_0,\tau_T\circ\rho)$. An arrow
$h:M\to N$ is also an arrow $T_*M\to T_*N$, since $\tau_T\circ\rho_N\circ h=\tau_T\circ\rho_M$, so $T_*$ is a functor,
with inverse $(T^{-1})_*$. It keeps $d(M)$ and the first-order deficit (\bgref{C.3}) and the gauge group (\bgref{D.21}),
and moves $S_1$ to $TS_1$. It does not keep rank, which depends on the matrix shape of $\Theta$ that $T$ need not
respect. Sums of $k$ Kronecker products are $(\mathcal R^{-1})_*\mathrm{LoRA}_k$, LoRA on the matrix rearranged by Van
Loan and Pitsianis and carried back (\bgref{B.11}), and HiRA \citep{huang2025hira} is $(D_{W_0})_*\mathrm{LoRA}$ with
$D_{W_0}(X)=W_0\odot X$, when $W_0$ has no zero entries. The atlas records $M$ and $T_*M$ as one object in two frames,
with no arrow between them.

\fbAusedin{\thmref{Proposition}{II.13}, Exposé~\ref{sec:cut}, rule T in \S\ref{sec:atlas-rules},
\S\ref{sec:rosetta-reverse}.}

\begin{bgentry}{A.18}{lift; lifting problem; obstruction}
Given arrows $p:E\to B$ and $f:A\to B$, a \textbf{lift} of $f$ along $p$ is an arrow $g:A\to E$ with $p\circ g=f$. A
\textbf{lifting problem} asks whether a lift exists, and an \textbf{obstruction} is an invariant whose value shows that
none does. A morphism $h:M\to N$ of $\Meth$ is exactly a pointed smooth lift of $\rho_M$ along $\rho_N$.
\end{bgentry}

In $\Set$ a lift exists once $f(A)\subseteq p(E)$, by choosing preimages, so smoothness is the real constraint. For
$\rho_M(t)=t$ and $\rho_N(u)=u^3$ on $\R$ the images agree, and the only lift is the cube root, which is not
differentiable at $0$. A pointed smooth lift forces $S_1(M)\subseteq S_1(N)$ by the chain rule, so $S_1$ is a
first-order obstruction, and here $S_1(M)=\R\not\subseteq S_1(N)=0$. For $t^3$ against $u^5$ both $S_1$ vanish, yet the
only lift $t^{3/5}$ is not $C^1$, so that obstruction is invisible to $S_1$. Merging a method on an extension is again a
lifting problem, to find $\mu$ with $\Phi_f\circ\mu=\Phi_{\tilde f}\circ\rho$, a lift along the realisation $\Phi_f$
(\bgref{B.16}, \thmref{Definition}{VI.1}).

\fbAusedin{\thmref{Observation}{I.4}(c), \S\ref{sec:slice-interval}, \thmref{Observation}{I.6},
\thmref{Definition}{VI.1}, \S\ref{sec:merge-extensions}, Rosetta entry 5 (\S\ref{sec:rosetta-six}).}

\begin{bgentry}{A.19}{base change; pushforward $\phi_!$; the relative point of view}
A map of bases $\phi:S'\to S$ induces the \textbf{pushforward} $\phi_!:\mathcal C/S'\to\mathcal C/S$,
$(X,x)\mapsto(X,\phi\circ x)$, by postcomposition. When every object over $S$ has a pullback along $\phi$, it also
induces the \textbf{pullback} $\phi^*:\mathcal C/S\to\mathcal C/S'$, $(Y,y)\mapsto(Y\times_SS',\text{projection})$, and
then $\phi_!\dashv\phi^*$. In algebraic geometry, base change usually means $\phi^*$.
\end{bgentry}

\thmref{Definition}{V.1}(a) uses $\phi_!$ on pointed slices. A function-preserving pointed reparametrisation
$\phi:(\Theta',\theta_0')\to(\Theta,\theta_0)$, with $f\circ(\phi\times X)=f'$, gives
$\phi_!:\Meth(\Theta',\theta_0')\to\Meth(\Theta,\theta_0)$, $\rho\mapsto\phi\circ\rho$, for instance for a permutation
of hidden units. Whether $\phi_!$ carries LoRA to LoRA is the covariance question of \thmref{Observation}{II.15}. A
quantiser changes the function and induces no such functor, so QLoRA \citep{dettmers2023qlora} and ReLoRA
\citep{lialin2023relora} only borrow the name, and Exposé~\ref{sec:base} computes pointing defects and reachable sets for
them.

Grothendieck's relative point of view replaces a space $X$ by a morphism $X\to S$, an object over a base. Properties then
belong to the morphism, and the useful ones are those that survive pulling back along every change of base $S'\to S$. The
fibre over a point $s$ is itself the pullback along $s:\mathrm{pt}\to S$. The paper takes this stance in a modest form. A
method is a morphism over the base point $\theta_0$, and quantisation, restarts and merging are read as moves of that
base point.

\fbAusedin{Exposés~\ref{sec:thesis}--\ref{sec:slice}, Exposé~\ref{sec:thesis} (where the category theory earns its
keep), \thmref{Definition}{V.1}, Exposé~\ref{sec:base}, Rosetta entry 28 (\S\ref{sec:rosetta-six}), the seven borrowed
words of Exposé~\ref{sec:analogy}. \textsf{Reading:} \citet{xvakil2025rising}, on fibred products and base change.}

\begin{bgentry}{A.20}{constructions the paper names and leaves unused}
Coends, monads, sheaves, stacks, schemes, Grothendieck's descent theory (descent data, effective descent morphisms), the
Grothendieck construction $\int\Meth\to\mathbf{PNet}$ and coKleisli 2-cells appear in the paper only in lists of
constructions it leaves out (Exposé~\ref{sec:thesis}, where the category theory earns its keep, and
Exposé~\ref{sec:analogy}, on what the borrowed words mean). No result depends on them. The elementary descent of
\bgref{D.22}, factoring a map that is constant on orbits through the orbit space, is a different and much simpler notion,
and several results use it.
\end{bgentry}

For orientation only, the Grothendieck construction would assemble the slices of all pretrained models,
$\Meth(\Theta,\theta_0)$ for each $(\Theta,\theta_0)$, into one category of pairs (network, method), with a projection to the category $\mathbf{PNet}$ whose objects are pointed networks and whose arrows are
function-preserving reparametrisations. A monad is an endofunctor $T$ with natural transformations
$\id\Rightarrow T$ and $TT\Rightarrow T$ satisfying unit and associativity laws, and a sheaf assigns data to open sets so
that compatible local data glue to unique global data.

\fbAusedin{Exposé~\ref{sec:thesis} (where the category theory earns its keep), the seven borrowed words and what the
borrowed words mean in Exposé~\ref{sec:analogy}.}

% ======================================================================================================================
\subsection*{Part B. Monoidal categories, string diagrams and lenses}
\phantomsection\label{app:bg-b}\addcontentsline{toc}{subsection}{Part B. Monoidal categories, string diagrams and lenses}

Part A compared methods through the arrows between them. Part B covers the algebra behind drawing a method as boxes and
wires, reading a rank bound off a cut of the drawing, and treating backpropagation as a pair of maps. Entries \bgref{B.1}
to \bgref{B.6} set up monoidal categories and their diagrams, \bgref{B.7} to \bgref{B.13} turn them into rank and
dimension counts for matrices built from small factors, and \bgref{B.14} to \bgref{B.17} cover sums of methods,
parametrised maps and lenses.

\subsubsection*{Monoidal categories}

\begin{bgentry}{B.1}{monoidal and symmetric monoidal categories}
A \textbf{monoidal category} is a category $\mathcal C$ with a functor $\otimes:\mathcal C\times\mathcal C\to\mathcal C$,
a unit object $I$, and natural isomorphisms $(A\otimes B)\otimes C\cong A\otimes(B\otimes C)$ and
$I\otimes A\cong A\cong A\otimes I$, the \textbf{coherence isomorphisms}, which satisfy the pentagon and triangle
equations. Mac Lane's coherence theorem then makes every diagram built from them commute, so brackets can be dropped. The
category is \textbf{symmetric} if it also has natural isomorphisms $\sigma_{A,B}:A\otimes B\to B\otimes A$ with
$\sigma_{B,A}\sigma_{A,B}=\id$ that are compatible with the associator (the hexagon equation). Examples are
$(\FinVect,\otimes,\R)$, $(\FinVect,\oplus,0)$ and $(\Diff,\times,\mathrm{pt})$.
\end{bgentry}

On methods, $M\oplus N$ runs the two parameter spaces side by side and adds their updates (\bgref{B.14}). This makes
$\Meth(\Theta,\theta_0)$ symmetric monoidal with the frozen model as unit. Its coherence isomorphisms are those of
$\times$ on pointed manifolds, and they lie over $\Theta$ because addition is associative and commutative.

\fbAusedin{\thmref{Observation}{I.7}, \S\ref{sec:slice-sums}, Exposé~\ref{sec:thesis} (where the category theory earns
its keep), \thmref{Theorem}{II.10}, \thmref{Theorem}{VI.3}. \textsf{Reading:} \citet[Ch.~VII and
Ch.~XI]{xmaclane1998categories}.}

\begin{bgentry}{B.2}{direct sum and tensor product in $\FinVect$}
The category $\FinVect$ has finite-dimensional real vector spaces as objects and linear maps as morphisms, and two
monoidal products.
The \textbf{direct sum} $V\oplus W$ is a \textbf{biproduct}, both the categorical product and the coproduct, with
$\dim(V\oplus W)=\dim V+\dim W$; on matrices $f\oplus g$ is block-diagonal. The \textbf{tensor product} $V\otimes W$ has
$\dim(V\otimes W)=\dim V\cdot\dim W$ and is not a categorical product; on matrices $f\otimes g$ is the Kronecker product
(\bgref{B.10}). Each bilinear map $V\times W\to U$ factors through a unique linear map $V\otimes W\to U$.
\end{bgentry}

Two boxes drawn side by side therefore mean a block-diagonal map in the activation diagrams of Exposé~\ref{sec:merge},
whose monoidal product is $\times=\oplus$, and a Kronecker product in a tensor network. Dimensions multiply along
$\otimes$, which is why the weight of a cut in \thmref{Theorem}{II.10} is a product of wire dimensions. LoRA's residual
branch $W_0+BA$ uses $\oplus$, since it sends one input to two branches and adds their outputs
(Figure~\ref{fig:bg-lora}).

\fbAusedin{\thmref{Theorem}{II.10}, Exposé~\ref{sec:cut}, \thmref{Theorem}{VI.3}, \thmref{Observation}{I.7}.}

\begin{bgentry}{B.3}{string diagrams}
A \textbf{string diagram} draws a morphism $f:A\to B$ of a monoidal category as a box with an input wire $A$ and an
output wire $B$. Composition $g\circ f$ joins boxes in series, $f\otimes g$ places them side by side, an identity is a
bare wire, the unit $I$ is drawn as no wire, and in a symmetric category $\sigma$ is a crossing. \citet{xjoyal1991geometry}
proved that the pictures are sound and complete. In a monoidal category, diagrams related by a planar deformation denote
the same morphism; in a symmetric one, diagrams with the same boxes joined in the same way do. Laws such as
$(g\otimes g')\circ(f\otimes f')=(g\circ f)\otimes(g'\circ f')$ can then be read off the picture.
\end{bgentry}

The plug of Exposé~\ref{sec:cut} (\S\ref{sec:cut-plug}) is a string diagram. LoRA is a box $A:\R^n\to\R^r$ followed by
$B:\R^r\to\R^m$, drawn beside $W_0$ (Figure~\ref{fig:bg-lora}, read left to right). The merge calculus moves an edit
along wires by equalities of maps, so a successful slide proves that the merged network computes the same function.

\fbAusedin{Exposé~\ref{sec:cut}, \thmref{Theorem}{II.10}, \S\ref{sec:cut-plug}, \S\ref{sec:merge-calculus} and the
analogy on string diagrams there, \thmref{Theorem}{VI.3}. \textsf{Reading:} \citet{xselinger2011survey}.}

\begin{bgentry}{B.4}{cartesian categories; copy and discard}
A monoidal category is \textbf{cartesian} when $\otimes$ is the categorical product $\times$ and $I$ is a terminal object
$1$. Every object then has a \textbf{copy} map $\Delta_X:X\to X\times X$, $x\mapsto(x,x)$, and a \textbf{discard} map
${!}_X:X\to1$, and copy is natural, $(f\times f)\circ\Delta_X=\Delta_Y\circ f$ for every $f:X\to Y$. Conversely, by a
theorem of \citet{xfox1976coalgebras}, a symmetric monoidal category is cartesian exactly when every object carries a
commutative comonoid, a copy and a discard map obeying coassociativity, counit and commutativity laws (\bgref{B.6}),
compatible with $\otimes$ and natural in all morphisms. $\Set$, $\Diff$ and $(\FinVect,\oplus)$ are cartesian.
$(\FinVect,\otimes)$ is not, since $x\mapsto x\otimes x$ is not linear.
\end{bgentry}

Naturality of copy is rule R7 of the merge calculus. A linear edit $E$ on a wire read by $W_Q$, $W_K$ and $W_V$ equals
$E$ applied on each of the three branches, so it is absorbed as $W_QE$, $W_KE$, $W_VE$. The paper builds $\Para(\mathcal C)$
(\bgref{B.16}) over a cartesian $\mathcal C$. The linear copy of \bgref{B.6} copies basis vectors only and is a different
map.

\fbAusedin{\thmref{Definition}{I.1}, \S\ref{sec:merge-calculus} and the analogy on string diagrams there,
\thmref{Theorem}{VI.3}.}

\begin{bgentry}{B.5}{compact closed categories; cups, caps and wire-bending}
A symmetric monoidal category is \textbf{compact closed} if every object $A$ has a dual $A^*$ with a \textbf{cup}
$\eta_A:I\to A^*\otimes A$ and a \textbf{cap} $\varepsilon_A:A\otimes A^*\to I$ satisfying the snake equations
\begin{equation*}
  (\varepsilon_A\otimes\id_A)(\id_A\otimes\eta_A)=\id_A,\qquad(\id_{A^*}\otimes\varepsilon_A)(\eta_A\otimes\id_{A^*})=\id_{A^*},
\end{equation*}
with the unitors $I\otimes A\cong A\cong A\otimes I$ of \bgref{B.1} suppressed. Drawn as bent wires, cups and caps turn
an input wire into an output wire and back. In $\FinVect$ the cap is evaluation, the cup is
$1\mapsto\sum_ie_i^*\otimes e_i$, and bending gives $\mathrm{Hom}(V,U)\cong U\otimes V^*$. With a fixed basis, a box
slid around a cup or cap comes out as its transpose.
\end{bgentry}

Bending the input wire of an $m\times n$ matrix gives its vectorisation in $\R^m\otimes\R^n$, and the Van
Loan--Pitsianis rearrangement (\bgref{B.11}) is a composite of bends and crossings. Rule R5 of the merge calculus,
$\langle Mq,k\rangle=\langle q,M\T k\rangle$, slides $M$ around the cap that computes an attention score. The cut-rank
bound uses only $\otimes$ and crossings.

\fbAusedin{\thmref{Theorem}{II.11}, \thmref{Theorem}{VI.3}, Exposé~\ref{sec:cut}. \textsf{Reading:}
\citet{xcoecke2017picturing}.}

\begin{bgentry}{B.6}{Frobenius algebras and spiders}
A \textbf{Frobenius algebra} on an object $V$ is a monoid $\mu:V\otimes V\to V$, $\eta:I\to V$ and a comonoid
$\delta:V\to V\otimes V$, $\varepsilon:V\to I$ obeying the \textbf{Frobenius law}
\begin{equation*}
  (\mu\otimes\id)(\id\otimes\delta)=\delta\mu=(\id\otimes\mu)(\delta\otimes\id).
\end{equation*}
A basis $e_1,\dots,e_n$ of $\R^n$ gives one, with \emph{copy} $\delta(e_i)=e_i\otimes e_i$, \emph{merge} (the
multiplication; unrelated to merging an adapter into the weights) $\mu(e_i\otimes e_j)=e_i$ if $i=j$ and $0$ otherwise,
$\eta=\sum_ie_i$ and $\varepsilon(e_i)=1$. It is commutative and special ($\mu\delta=\id$). For such algebras the
\textbf{spider theorem} says that every connected diagram of $\mu,\eta,\delta,\varepsilon$ with $k\ge1$ inputs and $l$
outputs equals one \textbf{spider}, which sends $e_i^{\otimes k}$ to $e_i^{\otimes l}$ and every other basis tensor to
$0$.
\end{bgentry}

Merge multiplies coordinates, $\mu(x\otimes y)=x\odot y$. So for matrices $X,Y:\R^n\to\R^m$ the Hadamard product is the
spider composite
\begin{equation*}
  X\odot Y=\mu_m\circ(X\otimes Y)\circ\delta_n.
\end{equation*}
LoHa \citep{hyeonwoo2021fedpara} is a copy spider on the input feeding $A_1$ and $A_2$, then $B_1$ and $B_2$, then a merge
spider on the output, and in \thmref{Theorem}{II.10} a spider split by the cut counts its dimension once. The copy and
$\odot$ nodes of Exposé~\ref{sec:merge} live in a cartesian category (\bgref{B.4}), where copy is $x\mapsto(x,x)$;
together they fail the Frobenius law, so there ``spider'' is only a name.

\fbAusedin{Exposé~\ref{sec:cut}, \thmref{Theorem}{II.10}, \thmref{Proposition}{II.12}, the analogy on string diagrams in
\S\ref{sec:merge-calculus}. \textsf{Reading:} \citet{xcoecke2013new}.}

\subsubsection*{Tensor networks and rank}

\begin{bgentry}{B.7}{tensor networks, contraction, bond dimension}
A \textbf{tensor network} $\Gamma$ is a string diagram in $(\FinVect,\otimes)$, possibly with spiders. Each box is a
tensor and each wire a finite-dimensional space. A wire joining two boxes means \textbf{contraction}, a sum over the index
the wire carries. The open wires (legs) are the inputs $V$ and outputs $U$, and contracting every internal wire gives the
value $\llbracket\Gamma\rrbracket:V\to U$. The dimension of an internal wire is its \textbf{bond dimension}. One value has
many presentations, and quantities read off a network, such as its cuts, depend on the presentation.
\end{bgentry}

LoRA's update is the two-box network $B\circ A$ with one bond of dimension $r$, and $\Delta W=BA$ contracts it,
$(BA)_{ij}=\sum_{k\le r}B_{ik}A_{kj}$. HiRA factors $W_0$ through a bond of dimension $\operatorname{rk}W_0$, while MoRA
\citep{jiang2024mora}, C3A \citep{chen2024parameter} and MPO adapters have no narrow cut between inputs and outputs
(Table~\ref{tab:cutbounds}, after \thmref{Theorem}{II.10}).

\fbAusedin{\thmref{Theorem}{II.10}, Exposé~\ref{sec:cut}, \S\ref{sec:cut-plug}, Exposé~\ref{sec:thesis} (where the
category theory earns its keep), the known results credited in Exposé~\ref{sec:related}. \textsf{Reading:}
\citet{xorus2014practical}.}

\begin{bgfact}{B.8}{the min-cut bound on rank}
Let $\Gamma$ be a tensor network with input legs $V$ and output legs $U$, with every wire attached to a box (put an
identity box on any bare strand). Split the boxes into an input side $\mathcal I$ and an output side $\mathcal O$. The
\textbf{weight} $w$ of the cut is the product of the dimensions of the wires joining the two sides and of the legs attached
on the wrong side (input legs at $\mathcal O$, output legs at $\mathcal I$); a spider with legs on both sides contributes
its dimension once. Then
\begin{equation*}
  \rk\llbracket\Gamma\rrbracket\le\min_{\text{cuts}}w.
\end{equation*}
For sums of networks, $\rk(X+Y)\le\rk X+\rk Y$.
\end{bgfact}

\paragraph{The idea of the proof}
Contract each side separately. The value becomes a map into the tensor product of the crossing wires and wrong-side legs
followed by a map out of it, and that space has dimension $w$.

\medskip\noindent
For LoRA the cut between $A$ and $B$ crosses one wire, so $\rk BA\le r$ (Figure~\ref{fig:bg-lora}). For LoHa the cut
with $A_1,A_2$ on the input side crosses two bonds and gives $r_1r_2$. The bound depends on the presentation. HRA's
product of $r$ reflections expands into $2^r-1$ terms of rank at most one, while the telescoping sum
\begin{equation*}
  R-I=\sum_{k=1}^r(H_k-I)H_{k+1}\cdots H_r
\end{equation*}
gives $r$.

\begin{figure}[tbp]
\centering
\begin{tikzpicture}[x=1cm,y=1cm,
  box/.style={rectangle, rounded corners=2pt, minimum height=8mm, line width=0.6pt, text=ink},
  frozen/.style={box, draw=inkgray, fill=headergray, minimum width=22mm},
  train/.style={box, draw=tide, fill=tide!16, line width=0.8pt, minimum width=11mm},
  wire/.style={draw=ink, line width=1.2pt},
  thinwire/.style={draw=ink, line width=0.5pt},
  lab/.style={font=\footnotesize, text=inkgray}]
  % input wire and fork
  \draw[wire] (0,0) -- (1.0,0);
  \fill[ink] (1.0,0) circle (1.9pt);
  % upper branch: frozen W_0
  \node[frozen] (W0) at (4.6,1.0) {$W_0$};
  \draw[wire] (1.0,0) .. controls (1.55,0) and (1.65,1.0) .. (2.25,1.0) -- (W0.west);
  \draw[wire] (W0.east) -- (6.95,1.0) .. controls (7.55,1.0) and (7.6,0) .. (7.86,0);
  % lower branch: A, thin wire r, B
  \node[train] (A) at (3.0,-1.0) {$A$};
  \node[train] (B) at (6.2,-1.0) {$B$};
  \draw[wire] (1.0,0) .. controls (1.55,0) and (1.65,-1.0) .. (2.25,-1.0) -- (A.west);
  \draw[thinwire] (A.east) -- (B.west);
  \draw[wire] (B.east) -- (6.95,-1.0) .. controls (7.55,-1.0) and (7.6,0) .. (7.86,0);
  % sum node and output
  \node[circle, draw=ink, fill=white, line width=0.7pt, minimum size=4.4mm, inner sep=0pt] (plus) at (8.08,0) {};
  \draw[ink, line width=0.7pt] (7.96,0) -- (8.20,0) (8.08,-0.12) -- (8.08,0.12);
  \draw[wire] (plus.east) -- (9.3,0);
  % dimension labels
  \node[lab, anchor=south] at (0.45,0.06) {$n$};
  \node[lab, anchor=south] at (8.85,0.06) {$m$};
  \node[lab, anchor=south] at (4.1,-0.98) {$r$};
  % the cut
  \draw[seal, line width=1pt, dash pattern=on 3pt off 2pt] (4.6,-0.35) -- (4.6,-1.75);
  \node[font=\footnotesize, text=seal, anchor=north] at (4.6,-1.78) {cut, weight $r$};
\end{tikzpicture}
\caption{LoRA as a string diagram, read left to right, with the scale $s$ absorbed into $B$. The fork sends the input to
  both branches and the $+$ node adds their outputs. These are the copy and sum maps of the biproduct $\oplus$
  (\bgref{B.2}), so the diagram says $W=W_0+BA$. The frozen $W_0$ is grey, and the trainable $A\in\R^{r\times n}$ and
  $B\in\R^{m\times r}$ are tinted. The dashed cut separates $A$ from $B$ and crosses one wire of dimension $r$, so
  $\rk(W-W_0)=\rk BA\le r$. For $W$ itself, subadditivity over the two branches gives only $\rk W\le\rk W_0+r$.}
\label{fig:bg-lora}
\end{figure}

\fbAusedin{\thmref{Theorem}{II.10}, Exposé~\ref{sec:cut}, \S\ref{sec:cut-plug}, \thmref{Proposition}{II.12}, the known
results credited in Exposé~\ref{sec:related}. \textsf{Reading:} \citet{kolda2009tensor}.}

\begin{bgentry}{B.9}{tensor trains, Tucker and MPO; gauge on bonds}
A \textbf{tensor train} (TT) writes a tensor with legs of dimensions $n_1,\dots,n_d$ as a chain of cores
$G_k\in\R^{r_{k-1}\times n_k\times r_k}$, $r_0=r_d=1$, joined by bonds of dimensions $r_1,\dots,r_{d-1}$. An
\textbf{MPO} is the same chain with one input and one output leg on each core. \textbf{Tucker} joins a core
$C\in\R^{r_1\times\cdots\times r_d}$ to factor matrices $U_k\in\R^{n_k\times r_k}$. Inserting $g^{-1}g=\id$,
$g\in\GL_{r_e}$, on a bond $e$ and absorbing each factor into its neighbouring box leaves the value unchanged, so the
gauge (\bgref{D.21}) contains $\prod_e\GL_{r_e}$. When every $r_e$ is the rank of the matching unfolding (minimal ranks),
the tensors of these ranks form a manifold of dimension $\sum(\text{core sizes})-\sum_er_e^2$.
\end{bgentry}

LoRA is the two-core case, $B$ and $A$ joined by one bond, with gauge $(B,A)\mapsto(Bg^{-1},gA)$ and, for
$r\le\min(m,n)$, image dimension $r(m+n)-r^2$. TT and Tucker adapters have $d(M)=\lvert M\rvert-\sum_er_e^2$ at minimal
ranks (\thmref{Proposition}{II.12}), and under gradient flow each bond carries a conserved matrix
(\S\ref{sec:fibres-charges}).

\fbAusedin{\thmref{Proposition}{II.12}, Exposé~\ref{sec:cut}, \S\ref{sec:fibres-charges}. \textsf{Reading:}
\citet{oseledets2011tensor}; \citet{xholtz2012manifolds}; \citet{xkoch2010dynamical}.}

\begin{bgentry}{B.10}{Kronecker product, vec, tensor product of subspaces}
For $C\in\R^{m_1\times n_1}$ and $D\in\R^{m_2\times n_2}$, the \textbf{Kronecker product}
$C\otimes D\in\R^{m_1m_2\times n_1n_2}$ is the block matrix $[c_{ij}D]$, the matrix of $C\otimes D$ in the product basis.
It obeys the mixed-product rule $(C\otimes D)(E\otimes F)=CE\otimes DF$, and $\rk(C\otimes D)=\rk C\cdot\rk D$. The map
$u\otimes v\mapsto uv\T$ identifies $\R^m\otimes\R^n$ with $\R^{m\times n}$. For subspaces $U\subseteq\R^m$ and
$V\subseteq\R^n$, $U\otimes V$ is then the space of matrices with column space in $U$ and row space in $V$, of dimension
$\dim U\cdot\dim V$.
\end{bgentry}

LoKr's \citep{yeh2023navigating} update is $C\otimes BA$ with $C$ ranging over all $m_1\times n_1$ matrices,
$B\in\R^{m_2\times r'}$ and $A\in\R^{r'\times n_2}$. For inner rank $r'\le\min(m_2,n_2)$ its rank
$\operatorname{rk}C\cdot\operatorname{rk}(BA)$ is at most $\min(m_1,n_1)\,r'$, and this is attained. So
$\Im\mathrm{LoKr}$ lies inside the image $W_0+\Mr$ of zero-init $\mathrm{LoRA}_r$ exactly when $r\ge\min(m_1,n_1)\,r'$
(\thmref{Theorem}{II.11}(e)). At a LoRA point $(B,A)$, the directions reached by moving one factor alone overlap in
\begin{equation*}
  \{XA\}\cap\{BY\}=\operatorname{col}B\otimes\operatorname{row}A,
\end{equation*}
of dimension $\operatorname{rk}A\cdot\operatorname{rk}B$ (\thmref{Proposition}{II.3}).

\fbAusedin{\thmref{Theorem}{II.11}, Exposé~\ref{sec:cut}, \thmref{Proposition}{II.3}, \thmref{Theorem}{II.2}, the
\hyperref[ex:loha]{LoHa\,/\,LoKr example} of Exposé~\ref{sec:examples}. \textsf{Reading:}
\citet[Ch.~4]{xhorn1991topics}.}

\begin{bgfact}{B.11}{Van Loan--Pitsianis rearrangement; Kronecker rank}
Let $m=m_1m_2$ and $n=n_1n_2$, and cut $\Delta\in\R^{m\times n}$ into an $m_1\times n_1$ grid of $m_2\times n_2$ blocks
$\Delta_{ij}$. The \textbf{rearrangement} $\mathcal R(\Delta)\in\R^{m_1n_1\times m_2n_2}$ has the rows
$\vect(\Delta_{ij})\T$, ordered with $i$ running fastest. Then
\begin{equation*}
  \mathcal R(C\otimes D)=\vect(C)\,\vect(D)\T,
\end{equation*}
and $\mathcal R$ is a linear bijection that permutes entries, hence a Frobenius isometry. So the \textbf{Kronecker rank}
$\operatorname{rk}_\otimes\Delta$, the least $k$ with $\Delta=\sum_{i\le k}C_i\otimes D_i$, equals
$\rk\mathcal R(\Delta)$, and the nearest sum of $k$ Kronecker products in Frobenius norm comes from the rank-$k$
truncated SVD of $\mathcal R(\Delta)$ (Van Loan and Pitsianis~\citeyear{loan1993approximation}).
\end{bgfact}

\paragraph{The idea of the proof}
Block $(i,j)$ of $C\otimes D$ is $c_{ij}D$, so its row in $\mathcal R$ is $c_{ij}\,\vect(D)\T$. Linearity extends this
to sums, and Eckart--Young applies to $\mathcal R(\Delta)$ because $\mathcal R$ preserves the Frobenius norm.

\medskip\noindent
In string-diagram terms $\mathcal R$ bends and crosses wires, regrouping the four legs of $\Delta$ from
$U_1U_2\mid V_1V_2$ to $U_1V_1\mid U_2V_2$. A sum of $k$ Kronecker products is therefore $\mathrm{LoRA}_k$ transported by
$\mathcal R^{-1}$ (\thmref{Proposition}{II.13}), with rank taken across another cut. For $N=s^2$, $I_N=I_s\otimes I_s$
has Kronecker rank $1$ and rank $N$. Figure~\ref{fig:cut} shows the singular values across the ordinary cut, this
Kronecker cut and a third one, for a random rank-one update, the identity and a Kronecker product $C\otimes D$. In this
paper a \emph{Kronecker sum} means a sum of Kronecker products; matrix analysis uses the same name for
$C\otimes I+I\otimes D$.

\fbAusedin{\thmref{Theorem}{II.11}, \thmref{Proposition}{II.13}, Figure~\ref{fig:cut}, \thmref{Prediction}{X.5}, the
known results and key references of Exposé~\ref{sec:related}. \textsf{Reading:} \citet{xvanloan2000ubiquitous}.}

\begin{bgentry}{B.12}{Schmidt and operator-Schmidt decompositions}
Every $u\in\R^{m_1}\otimes\R^{m_2}$ has a \textbf{Schmidt decomposition} $u=\sum_{i\le\rho}\sigma_ix_i\otimes y_i$ with
$\sigma_i>0$ and orthonormal families $(x_i)$, $(y_i)$. It is the SVD of the $m_1\times m_2$ reshape of $u$
($x\otimes y\mapsto xy\T$), and the \textbf{Schmidt rank} $\rho_u$ is the rank of that reshape, at most
$\min(m_1,m_2)$. The SVD of $\mathcal R(\Delta)$ likewise gives the \textbf{operator-Schmidt decomposition}
$\Delta=\sum_{k\le K}s_kC_k\otimes D_k$, with $s_k>0$, Frobenius-orthonormal families $(C_k)$ and $(D_k)$, and
$K=\operatorname{rk}_\otimes\Delta$.
\end{bgentry}

If $u=\sum_i\sigma_ix_i\otimes y_i$ and $v=\sum_j\tau_jz_j\otimes w_j$ have Schmidt ranks $\rho_u$ and $\rho_v$, the
mixed-product rule gives
\begin{equation*}
  uv\T=\sum_{i,j}\sigma_i\tau_j\,(x_iz_j\T)\otimes(y_iw_j\T),
\end{equation*}
an operator-Schmidt decomposition with $\rho_u\rho_v$ terms, so $\operatorname{rk}_\otimes(uv\T)=\rho_u\rho_v$. A generic
rank-one update, which $\mathrm{LoRA}_1$ reaches, thus needs $\kappa_1=\min(m_1,m_2)\min(n_1,n_2)$ Kronecker terms
(\thmref{Theorem}{II.11}(c)).

\fbAusedin{\thmref{Theorem}{II.11}, Figure~\ref{fig:cut}, \thmref{Prediction}{X.5}, the known results credited in
Exposé~\ref{sec:related}. \textsf{Reading:} \citet[\S2.5]{xnielsen2000quantum}.}

\begin{bgfact}{B.13}{rank of a Hadamard product; face-splitting product}
For $X,Y\in\R^{m\times n}$, $\rk(X\odot Y)\le\rk X\cdot\rk Y$. Explicitly, if $X=B_1A_1$ and $Y=B_2A_2$ with inner
dimensions $r_1$ and $r_2$, then
\begin{equation*}
  X\odot Y=(B_1\bullet B_2)\,(A_1\T\bullet A_2\T)\T,
\end{equation*}
where the \textbf{face-splitting product} $B_1\bullet B_2\in\R^{m\times r_1r_2}$ has as row $i$ the Kronecker product of
row $i$ of $B_1$ and row $i$ of $B_2$. In particular
\begin{equation*}
  W_0\odot(BA)=\sum_{k\le r}\diag(b_k)\,W_0\,\diag(a_k),
\end{equation*}
summed over the columns $b_k$ of $B$ and the rows $a_k$ of $A$, has rank at most $r\rk W_0$
\citep[Ch.~5]{xhorn1991topics}.
\end{bgfact}

\paragraph{The idea of the proof}
$X\odot Y$ is the submatrix of $X\otimes Y$ on the rows $(i,i)$ and the columns $(j,j)$. In pictures it is the spider
composite $\mu\circ(X\otimes Y)\circ\delta$ of \bgref{B.6}, which factors through the two inner wires.

\medskip\noindent
$\mathrm{LoHa}_{r_1,r_2}$ therefore has rank at most $r_1r_2$, and $h=\big(B_1\bullet B_2,(A_1\T\bullet A_2\T)\T\big)$ is
the atlas arrow $\mathrm{LoHa}_{r_1,r_2}\to\mathrm{LoRA}_{r_1r_2}$. HiRA's update $W_0\odot BA$ has rank at most
$r\rk W_0$, a bound that allows full rank when $W_0$ has full rank (\thmref{Proposition}{II.9}(a)).

\fbAusedin{\thmref{Proposition}{II.12}, \thmref{Theorem}{II.10}, \thmref{Proposition}{II.9}, \S\ref{sec:atlas-arrows},
Rosetta entries 6 and 11 (\S\ref{sec:rosetta-six}).}

\subsubsection*{Sums, Para and lenses}

\begin{bgentry}{B.14}{sums of methods; Minkowski sums, differences and translates}
For subsets $S,T$ of a vector space, the \textbf{Minkowski sum} is $S+T=\{s+t:s\in S,\ t\in T\}$, the
\textbf{Minkowski difference} is $S-T=\{s-t:s\in S,\ t\in T\}$, a \textbf{translate} of $S$ is $v+S=\{v+s:s\in S\}$,
and $C^{+K}$ is the sum of $K$ copies of $C$. A sum is the image of $S\times T$ under addition, so for semialgebraic sets
$\dim(S+T)\le\dim S+\dim T$ (\bgref{C.18}), and likewise for differences. If $0\in C$, the additive monoid
(\bgref{D.10}) generated by $C$ is $\operatorname{Mon}(C)=\bigcup_KC^{+K}$, and if moreover $\R C=C$ it is
$\spanop C$. Writing $\rho_M=\theta_0+\delta_M$, the sum of methods
\begin{equation*}
  M\oplus N=\big(Q_M\times Q_N,\,(q_0,q_0'),\,\theta_0+\delta_M+\delta_N\big)
\end{equation*}
has image $\theta_0+(\Im\delta_M+\Im\delta_N)$, and $\oplus$ is a symmetric monoidal structure with the frozen model as
unit (\bgref{B.1}).
\end{bgentry}

Concatenation $([B_1\ B_2],[A_1;A_2])$ is an isomorphism $\mathrm{LoRA}_r\oplus\mathrm{LoRA}_s\cong\mathrm{LoRA}_{r+s}$
of unpointed objects, with the $A$-blocks rescaled to the target scale when the scales differ, as under $\alpha/r$
scaling (\thmref{Observation}{I.7}). It is also a pointed isomorphism onto $\mathrm{LoRA}_{r+s}$ pointed at the
concatenated initialisation. RoSA \citep{nikdan2024rosa} and GraLoRA \citep{jung2025gralora} are $\oplus$-words. The
categorical product of \thmref{Observation}{I.4}, when it exists, has image $\Im M\cap\Im N$ (\bgref{A.10}). Merging and
restarting $K$ times with a fixed shape $C$ reaches $\theta_0+C^{+K}$. Since
$\Mr^{+K}=\mathcal M_{\le\min(Kr,m,n)}$, ReLoRA reaches every matrix after $\lceil\min(m,n)/r\rceil$ cycles
(\thmref{Theorem}{V.3}).

Re-pointing an additive method $\rho=\theta_b+\delta$ with shape $C=\Im\delta$ at a parameter $q$, so that
$\theta_b=\theta_0-\delta(q)$, gives the translate $\theta_0+(C-\delta(q))$ as image, and all exact pointings together
reach $\theta_0+(C-C)$. For LoRA, $\Mr$ is closed under negation and $\Mr+\Mr=\mathcal M_{\le\min(2r,m,n)}$ by
splitting an SVD (\bgref{C.12}), so $C-C=\mathcal M_{\le\min(2r,m,n)}$, which is why no exact pointing reaches an update
of rank above $2r$ (\thmref{Theorem}{II.1}).

\fbAusedin{\thmref{Observation}{I.7}, \S\ref{sec:slice-sums}, \thmref{Theorem}{II.1}, \thmref{Proposition}{II.13},
\thmref{Theorem}{V.3}, \S\ref{sec:base-rebase}, \S\ref{sec:rosetta-reverse}.}

\begin{bgentry}{B.15}{strong monoidal functors; $(\N,+)$ and the permutation groupoid}
A \textbf{strong monoidal functor} $F:(\mathcal C,\otimes,I)\to(\mathcal D,\otimes,J)$ is a functor with isomorphisms
$\phi_{A,B}:FA\otimes FB\to F(A\otimes B)$, natural in $A$ and $B$, and $\phi_0:J\to FI$, compatible with associators and
unitors. It is \textbf{symmetric} if also $F(\sigma_{A,B})\circ\phi_{A,B}=\phi_{B,A}\circ\sigma_{FA,FB}$. The
\textbf{discrete monoidal category} $(\N,+)$ has the natural numbers as objects, only identity morphisms, $r\otimes s=r+s$
and identity symmetry. The \textbf{permutation groupoid} $\mathbb P$ has the same objects, the symmetric group $S_r$ as
the automorphisms of $r$ and no other morphisms, with $\otimes$ given by $+$ and block sums of permutations, and the block
transposition in $S_{r+s}$ as symmetry.
\end{bgentry}

Concatenation (\bgref{B.14}) makes $r\mapsto\mathrm{LoRA}_r$ strong monoidal from $(\N,+)$ into unpointed methods. It is
not symmetric there, because $\phi_{s,r}\circ\sigma$ lists the blocks in the other order. On $\mathbb P$ the block
transposition acts on $\mathrm{LoRA}_{r+s}$ by permuting inner channels, $(B,A)\mapsto(BP^{-1},PA)$, and the symmetry
square commutes. Pointed objects behave differently. If $F(r)$ and $F(2r)$ are standard $\mathrm T_A$ pointings
($B_0=0$ and $A_0$ of full row rank, \bgref{C.19}), which needs $2r<\min(m,n)$, then $F(r)\oplus F(r)\not\cong F(2r)$ as
pointed objects, since $\dim S_1$ is $mr$ against $2mr$ (\thmref{Observation}{I.7}).

\fbAusedin{\thmref{Observation}{I.7}, the referee record in Exposé~\ref{sec:analogy}. \textsf{Reading:}
\citet[Ch.~XI]{xmaclane1998categories}.}

\begin{bgentry}{B.16}{bicategories and the Para construction}
A \textbf{bicategory} has objects, \textbf{1-cells} between objects and \textbf{2-cells} between parallel 1-cells.
2-cells compose vertically and horizontally, and composition of 1-cells is associative and unital up to coherent
invertible 2-cells. For a symmetric monoidal category $\mathcal C$, here always cartesian, the bicategory
$\Para(\mathcal C)$ has the objects of $\mathcal C$. A 1-cell $X\to Y$ is a pair $(P,\,f:P\times X\to Y)$, and
$(P',g):Y\to Z$ composes with it to $\big(P'\times P,\ (p',p,x)\mapsto g(p',f(p,x))\big)$. A 2-cell
$(P,f)\Rightarrow(Q,g)$ is a \textbf{reparametrisation} $r:Q\to P$ with $g=f\circ(r\times\id_X)$. A \textbf{2-functor}
(pseudofunctor) between bicategories sends objects, 1-cells and 2-cells to objects, 1-cells and 2-cells, and preserves
both compositions up to coherent invertible 2-cells.
\end{bgentry}

A linear layer is the 1-cell $\big(\R^{m\times n},\,(W,x)\mapsto Wx\big):\R^n\to\R^m$, and stacking a second layer
$\R^m\to\R^k$ composes 1-cells, with parameter space $\R^{k\times m}\times\R^{m\times n}$. A pretrained model is a
pointed 1-cell $F=(\Theta,\theta_0,f)$ of $\Para(\Diff)$, and its realisation $\Phi_f:\Theta\to C^\infty(X,Y)$,
$\theta\mapsto f(\theta,-)$, records the function at each weight. LoRA is a 2-cell
$F\Rightarrow(Q,f\circ(\rho\times\id_X))$, the network with $(B,A)$ in its parameter slot, and a simulation is a 2-cell
between two such cells (\thmref{Remark}{I.3}). An adapter extension is a 2-cell into $F$, and merging it is a lift
through $\Phi_f$ (\thmref{Definition}{VI.1}). Composition is associative only up to
$(P''\times P')\times P\cong P''\times(P'\times P)$, hence a bicategory; over a strict monoidal category it is a
2-category.

\fbAusedin{\thmref{Definition}{I.1}, \S\ref{sec:slice-para}, \thmref{Remark}{I.3}, \thmref{Definition}{VI.1},
\S\ref{sec:lens-functor}, the key references of Exposé~\ref{sec:related}. \textsf{Reading:}
\citet{xleinster1998basic}.}

\begin{bgentry}{B.17}{lenses, reverse derivatives and $\Para(\mathsf R)$}
In a cartesian category, a \textbf{lens} $(A,A')\to(B,B')$ is a pair of maps, a forward part $f:A\to B$ and a backward
part $f^\sharp:A\times B'\to A'$. Lenses compose by
\begin{equation*}
  (g,g^\sharp)\circ(f,f^\sharp)=\big(g\circ f,\ (a,c')\mapsto f^\sharp(a,g^\sharp(f(a),c'))\big)
\end{equation*}
and form the category $\Lens(\mathcal C)$. For smooth maps between open subsets of Euclidean spaces, the \textbf{reverse
derivative} is
\begin{equation*}
  \mathsf R[f](x,y)=J_f(x)\T y.
\end{equation*}
The chain rule $J_{g\circ f}(x)\T=J_f(x)\T J_g(f(x))\T$ says that
$f\mapsto(f,\mathsf R[f])$ is a functor into $\Lens$. Applied to parametrised maps it gives the 2-functor
$\Para(\mathsf R):\Para(\mathcal C)\to\Para(\Lens(\mathcal C))$ of \citet{cruttwell2021categorical}, which is
backpropagation.
\end{bgentry}

LoRA's forward part is $\rho(B,A)=W_0+sBA$, and its backward part pulls the weight gradient $G$ back to
\begin{equation*}
  \mathsf R[\rho]\big((B,A),G\big)=\big(sGA\T,\,sB\T G\big).
\end{equation*}
GaLore \citep{zhao2024galore} keeps the forward part $\id_\Theta$ and uses the compression $G\mapsto P\T G$ with anchor
$P$, a lens whose backward part is not the reverse derivative of its forward part (\thmref{Definition}{IV.1}). For linear
$\rho(C)=W_k+\mathcal LC$, $\mathsf R[\rho](C,G)=\mathcal L^*G$ does not depend on $C$, which is why the two readings
agree at degree one (\thmref{Proposition}{IV.2}).

\fbAusedin{Exposé~\ref{sec:lens}, \S\ref{sec:lens-functor}, \thmref{Definition}{IV.1}, \thmref{Proposition}{IV.2},
Rosetta entry 25 (\S\ref{sec:rosetta-six}), Exposé~\ref{sec:thesis} (where the category theory earns its keep), the
seven borrowed words of Exposé~\ref{sec:analogy}. \textsf{Reading:} \citet{cruttwell2021categorical}.}
